% FY27 ONR YIP informal brief topic description. Compile using pdflatex.
% Scientific framing tailored for goolsby. Program Officer line can be updated after consultation.
\documentclass[11pt,letterpaper]{article}
\usepackage[letterpaper,top=0.60in,bottom=0.58in,left=0.72in,right=0.72in]{geometry}
\usepackage{mathptmx}
\usepackage[T1]{fontenc}
\usepackage{microtype}
\usepackage{enumitem}
\usepackage{xcolor}
\usepackage[colorlinks=true,citecolor=blue,linkcolor=black,urlcolor=blue]{hyperref}
\pagestyle{empty}
\setlength{\parindent}{0pt}
\setlength{\parskip}{3pt}
\setlist[itemize]{leftmargin=1.13em,topsep=2pt,itemsep=2pt,parsep=0pt}
\emergencystretch=1.5em
\newcommand{\head}[1]{\textbf{#1}\enspace}
\newcommand{\POline}{Dr. Rebecca Goolsby, Code 341 --- Analytics, Modeling and Simulation (Local and Regional)}
\newcommand{\FocusArea}{Human and Biological Systems --- Warfighter Decision Superiority and Readiness}
\begin{document}
\begin{center}
{\fontsize{14}{16}\selectfont\bfseries Reasoning over Contested Evidence for Regional Situation Awareness\par}
\vspace{2pt}
{\small Brief topic description for Program Officer discussion --- ONR FY27 YIP (N0001426SF004)}
\end{center}
\vspace{-2pt}
{\small \textbf{Xinya Du}, Assistant Professor of Computer Science, The University of Texas at Dallas\par
\textbf{Email:} \href{mailto:xinya.du@utdallas.edu}{xinya.du@utdallas.edu} \quad \textbf{Website:} \href{https://xinyadu.github.io}{https://xinyadu.github.io}\par
\textbf{Program Officer:} \POline\par
\textbf{Proposed ONR S\&T Focus Area:} \FocusArea\par}
\vspace{2pt}\hrule\vspace{5pt}
\head{Problem.} Regional analysts synthesize news, public statements, social posts, and field reports that may be delayed, contradictory, or deliberately misleading. After a typhoon, dozens of posts may report contaminated drinking water, although they all repeat one unverified claim while a later field test finds no contamination. Conventional retrieval-augmented generation (RAG) can mistake repetition for independent corroboration and overlook corrections. The research challenge is to reason about what happened while also tracking how claims spread, changed, and acquired apparent credibility.
\par\head{Central hypothesis.} We hypothesize that coupling event-level reasoning with the lineage of reported claims can distinguish genuine corroboration from information amplification. Rather than detecting misinformation only after generating a summary, agents should retrieve and reconcile evidence according to source independence, temporal changes, and competing explanations, explicitly preserving uncertainty when no independent check is available.
\par\head{Research plan.}
\begin{itemize}
\item \textbf{Reconstruct events and information pathways.} Extract event claims and their actors, locations, times, and sources from heterogeneous reports. Link compatible and conflicting claims to evolving event hypotheses, while recording corrections, retractions, and likely quotation or copying relationships. Model uncertainty about these relationships so repeated reports do not automatically become independent confirmation or proof of intentional coordination.
\item \textbf{Develop lineage-aware retrieval and reasoning.} Rank evidence by its relevance, independence, timeliness, and ability to resolve competing event hypotheses, rather than by document similarity alone. Train retrieval and synthesis models with supervision on claim support, source dependence, and revisions. Produce grounded updates that distinguish corroborated events from disputed or unverified claims, and trigger targeted searches for independent observations when an apparent consensus rests on a single information pathway.
\item \textbf{Learn from counterfactual information environments.} Build a graph-based simulator linking underlying events, reporting actors, and information propagation. Generate controlled streams with delays, selective omissions, source copying, corrections, and coordinated claim amplification. Construct counterfactual cases in which message volume changes but independent evidence does not; use them to train evidence-selection policies and model how alternative information interventions change the inferred regional situation.
\end{itemize}
\head{Evaluation and expected outcomes.} Use \href{https://www.gdeltproject.org/}{\textcolor{blue}{GDELT}}, \href{https://trec.nist.gov/data/crisis.html}{\textcolor{blue}{TREC CrisisFACTS}}, and controlled report streams with known claim lineages and corrections. Compare against standard and self-refining RAG, focusing on support traceability, resistance to amplified false claims, correct handling of updates, and timely recognition of independently verified alerts.
\par\head{Naval relevance.} During humanitarian assistance and disaster relief, Navy and Marine Corps planners must distinguish verified local hazards from rumors amplified through online networks. Tracing the evidence behind evolving reports can improve regional awareness without suppressing legitimate early warnings; the proposed information-environment simulator also supports what-if analysis and mission rehearsal.
\par\head{Prior work and PI qualifications.} Our QA-based event extraction (\href{https://doi.org/10.18653/v1/2020.emnlp-main.49}{\textcolor{blue}{Du and Cardie, 2020}}) and DARPA KAIROS \href{https://doi.org/10.18653/v1/2022.naacl-demo.7}{\textcolor{blue}{RESIN-11}} work provide foundations for event representation across reports; FaithScore (\href{https://aclanthology.org/2024.findings-emnlp.290/}{\textcolor{blue}{Jing et al., 2024}}) and FGAIF (\href{https://openreview.net/forum?id=Qhfw5CUVd7}{\textcolor{blue}{Jing and Du, 2025}}) contribute experience with fine-grained grounding feedback. The PI holds a 2024 NSF CAREER award.

\end{document}
